\subsection{SISTEMAS \en{OFFLINE-FIRST}}

\en{Offline-first} é um paradigma de programação que especifica, dentre outros aspectos, o comportamento de aplicações em relação a falta de conectividade \cite{OfflineFirst}. Um bom exemplo desse paradigma são os \en{Progressive Web Apps (PWAs)}, para uma aplicação ser um \en{PWA} esta deve englobar conceitos que formulam o modo que foi construída; dentre eles destaca-se, para esse trabalho, a independência de conexão \cite{BjornHansen2018PWAUnifiedMobile}. A arquitetura define, também, o armazenamento local de dados como componente essencial para sua implementação, e com ele a sincronização e a resolução de conflitos. \en{Offline-first} assume a existência da necessidade de conexão com alguma rede, sendo essa efêmera, e especifica como aplicações devem funcionar nesse contexto \cite{Guda2025OfflineFirstAndroidChat}, um exemplo deste tipo de conexão são as redes \en{peer-to-peer}.  

A oferta variável de recursos da computação móvel sugere que aplicações para esse contexto devem ser adaptáveis. Historicamente, adaptação é a troca de um recurso por outro ou pela qualidade do serviço prestado, isso é: técnicas de redução, filtragem e transformação de dados para trafegarem na rede; o processo de adaptabilidade é disparado automaticamente após a leitura do contexto. Contexto é toda informação relevante para aplicação que define seu comportamento; desse modo, as informações de contexto podem ser categorizadas em: físicas, relativas ao \en{hardware}; lógicas, relativas ao \en{software}; e, sociais, relativas à comunidade de usuários e seu contexto \cite{AugustinYaminGeyerBarbosa2001RequisitosAplicacoesMoveisDistribuidas}. Essas considerações são essenciais para garantir o funcionamento de aplicações que, ao seguirem o paradigma \en{Offline-first}, continuem funcionando mesmo na ausência de conexão, requisito essencial para esses tipos de soluções. 
